% Generated by publication/build_sections.py; do not edit by hand.
\paragraph{Informal verdict.}
\textbf{A --- Complete solution}, with high confidence.

\paragraph{Lean coverage.}
Full canonical Q764 coverage is now ingested. Q764.Line.lineKCenter\_correct and lineKCenter\_work\_le connect the actual counted line run to the original Finset Real objective and prove <=100*(n*n+k*n) work; exhaustiveMetric\_correct and exhaustiveMetric\_work\_le give the actual exact finite-metric enumeration and its 24*(n choose k)*n*k+24*\(n+1\) bound; vertexCover\_iff\_oneTwo\_radius\_one, vertexCover\_polyReduces\_metricKCenter, and vertexCover\_inP\_of\_metricKCenter\_inP give the concrete 1--2 metric equivalence, polynomial materialization, and unconditional decision-solver transfer. No P != NP premise or unconditional impossibility theorem occurs.

\subsection{Audited informal answer}
\subsubsection{1. Faithful restatement}

Let \(1\le k\le n\).

\textbf{\(a\)} Given distinct real numbers
\[
X={x_1<x_2<\cdots <x_n},
\]
and \(F\subseteq X\), define
\[
\delta(X,F)=\max_{x\in X}\min_{y\in F}|x-y|.
\]
Find an efficient algorithm producing a \(k\)-element subset \(F\subseteq X\) minimizing \(\delta(X,F)\).

\textbf{\(b\)} Solve the analogous problem when \(X\) is an arbitrary \(n\)-point subset of a metric space, with (|x-y|) replaced by the metric distance. 

There are three minor ambiguities in the wording.

\begin{enumerate}
\item The statement only says \(k\le n\), but \(k\) must be a positive integer, since the minimum over an empty \(F\) is otherwise undefined.
\item The minimizing subset need not be unique, so ``the subset \(F\)'' should be read as ``a minimizing subset \(F\).''
\item ``Efficient'' is not formally defined. Below, arithmetic and comparisons of input distances are counted at unit cost. For rational inputs, the algorithms also have polynomial bit complexity.
\end{enumerate}

\medskip\hrule\medskip

\subsubsection{2. Final result}

\paragraph{Theorem}

\begin{enumerate}
\item \textbf{Ordered real line.} An optimal \(k\)-element set can be found exactly in
\end{enumerate}
\[
   O(n^2+kn)=O(n^2)
\]
   time and \(O(n^2+kn)=O(n^2)\) memory. The algorithm is a bottleneck dynamic program over consecutive blocks.

\begin{enumerate}
\item \textbf{Arbitrary finite metric.}
\end{enumerate}

\begin{itemize}
\item The problem is exactly the finite \textbf{discrete metric \(k\)-center problem}.
\item Its decision version is NP-complete, even when every nonzero distance is either \(1\) or \(2\). Consequently, unless \(P=NP\), there is no polynomial-time exact algorithm when \(k\) is part of the input.
\item An exact exhaustive algorithm takes
\end{itemize}
\[
     O\!\left(\binom nk nk\right)
\]
     time.
\begin{itemize}
\item There is an \(O(nk)\)-time polynomial algorithm returning a set \(F\) with
\end{itemize}
\[
     \delta(X,F)\le 2,\delta_{\mathrm{opt}}.
\]
     Unless \(P=NP\), no polynomial algorithm can guarantee a factor strictly smaller than \(2\).

\medskip\hrule\medskip

\subsection{Part \(a\): points on the real line}

\subsubsection{3. The optimal cost of one consecutive block}

For \(1\le i\le j\le n\), define
\[
c(i,j)
=
%
\min_{i\le \ell\le j}
\max_{i\le t\le j}|x_t-x_\ell|.
\]
Thus \(c(i,j)\) is the smallest radius with which one center, constrained to be one of
\[
x_i,x_{i+1},\ldots,x_j,
\]
can cover that whole block.

For a fixed center \(x_\ell\), the farthest point of the block is one of its two endpoints, so
\[
\max_{i\le t\le j}|x_t-x_\ell|
=
%
\max{x_\ell-x_i,;x_j-x_\ell}.
\]
Writing
\[
m_{ij}=\frac{x_i+x_j}{2},
\]
we have
\[
\max{x_\ell-x_i,x_j-x_\ell}
=
%
\frac{x_j-x_i}{2}
+
\left\lvert{}x_\ell-m_{ij}\right\rvert{}.
\]
Consequently:

\paragraph{Lemma 1}

An optimal one-center location for the block \(x_i,\ldots,x_j\) is any \(x_\ell\) nearest to the midpoint \((x_i+x_j)/2\), and
\[
c(i,j)
=
%
\frac{x_j-x_i}{2}
+
\min_{i\le \ell\le j}
\left\lvert{}x_\ell-\frac{x_i+x_j}{2}\right\rvert{}.
\tag{1}
\]

In particular, the center is constrained to lie in \(X\); this is not the continuous \(k\)-center problem.

\medskip\hrule\medskip

\subsubsection{4. Reduction to consecutive blocks}

The order on the real line provides the crucial structure.

\paragraph{Lemma 2: consecutive-cell lemma}

For every set of centers
\[
F={f_1<\cdots <f_s}\subseteq X,
\]
the points of \(X\) can be assigned to nearest centers so that the points assigned to each \(f_t\) form a consecutive block of \(X\).

\subparagraph{Proof}

Break ties consistently, for example in favor of the smaller center. The points at least as close to \(f_t\) as to \(f_{t+1}\) lie to the left of the midpoint
\[
\frac{f_t+f_{t+1}}2.
\]
Thus the nearest-center index is nondecreasing as one moves from left to right through \(X\). Its level sets are therefore consecutive intervals of indices.

Every center belongs to its own cell, since its distance to itself is zero, so all cells are nonempty. \(\square\)

It follows that an optimal solution is equivalent to partitioning \(X\) into at most \(k\) consecutive nonempty blocks and placing an optimal one-center solution in every block.

More precisely, if
\[
0=i_0<i_1<\cdots<i_s=n,
\]
the associated blocks are
\[
{x_{i_{t-1}+1},\ldots,x_{i_t}},
\qquad 1\le t\le s.
\]

\paragraph{Proposition 3: bottleneck-partition characterization}

The optimal value is
\[
R_k
=
%
\min_{\substack{1\le s\le k\\
0=i_0<i_1<\cdots<i_s=n}}
;
\max_{1\le t\le s}
c(i_{t-1}+1,i_t).
\tag{2}
\]

\subparagraph{Proof}

Given any set of at most \(k\) centers, Lemma 2 produces at most \(k\) consecutive cells. On a cell ([i,j]), its chosen center has radius at least \(c(i,j)\). Therefore its total radius is at least the right-hand side of \(2\).

Conversely, given a partition occurring on the right-hand side, choose in every block a center realizing \(c(i,j)\). Every point is then within the displayed maximum radius of its block center. \(\square\)

Using fewer than \(k\) centers causes no difficulty: arbitrary unused points of \(X\) can be added as extra centers, and doing so cannot increase the radius. Thus the minimum over ``at most \(k\)'' and the minimum over ``exactly \(k\)'' coincide.

\medskip\hrule\medskip

\subsubsection{5. Dynamic programming recurrence}

Let
\[
\rho_p(j)
\]
be the optimal radius for the prefix
\[
{x_1,\ldots,x_j}
\]
using at most \(p\) centers. Set
\[
\rho_p(0)=0,\qquad
\rho_0(j)=+\infty\quad(j\ge1).
\]

The last block of a partition of \(x_1,\ldots,x_j\) must have the form
\[
x_i,x_{i+1},\ldots,x_j
\]
for some \(1\le i\le j\). Hence:

\paragraph{Lemma 4: dynamic-programming recurrence}

For \(p,j\ge1\),
\[
\boxed{
\rho_p(j)
=
%
\min_{1\le i\le j}
\max\bigl\{\rho_{p-1}(i-1),c(i,j)\bigr\}.
}
\tag{3}
\]

\subparagraph{Proof}

If the last block starts at \(i\), the preceding prefix uses at most \(p-1\) centers and has optimal cost \(\rho_{p-1}(i-1)\); the final block costs \(c(i,j)\). The total bottleneck cost is their maximum.

Conversely, every partition counted by \(\rho_p(j)\) has a uniquely determined last-block starting index \(i\), so every possible solution appears in the minimum. \(\square\)

A direct evaluation of \(3\) would take \(O(kn^2)\) time. We now exploit monotonicity to reduce the dynamic-programming part to \(O(kn)\).

\medskip\hrule\medskip

\subsubsection{6. Monotonicity of the transition}

For fixed \(p\) and \(j\), put
\[
A_i=\rho_{p-1}(i-1),
\qquad
B_i=c(i,j).
\]

\paragraph{Lemma 5}

As \(i\) increases:

\begin{enumerate}
\item \(A_i\) is nondecreasing;
\item \(B_i\) is nonincreasing.
\end{enumerate}

Moreover, for fixed \(i\), \(c(i,j)\) is nondecreasing in \(j\).

\subparagraph{Proof}

The value \(A_i\) concerns a prefix ending at \(i-1\). Enlarging that prefix cannot decrease the required radius, so \(A_i\) is nondecreasing.

For \(B_i\), removing the leftmost point of a block cannot increase the optimum. If an optimal center for ([i,j]) is not \(x_i\), retain it. If it is \(x_i\), replace it by \(x_{i+1}\); every remaining point lies to the right, and is no farther from \(x_{i+1}\) than from \(x_i\). Thus
\[
c(i+1,j)\le c(i,j).
\]

Similarly, adding a point on the right cannot reduce the one-center cost. If an optimal center for ([i,j+1]) is one of \(x_i,\ldots,x_j\), restricting to the smaller block proves the assertion. If it is \(x_{j+1}\), then \(x_j\) is at least as good for all points \(x_i,\ldots,x_j\). Hence
\[
c(i,j)\le c(i,j+1).
\]
\(\square\)

Let
\[
t=t_p(j)=
\min{i\in{1,\ldots,j}: A_i\ge B_i},
\]
with the convention \(t=j+1\) if this set is empty.

For \(i<t\),
\[
\max(A_i,B_i)=B_i,
\]
which is nonincreasing in \(i\). For \(i\ge t\),
\[
\max(A_i,B_i)=A_i,
\]
which is nondecreasing in \(i\). Therefore:

\paragraph{Lemma 6: crossing lemma}

A minimizing index in \(3\) is always one of
\[
t_p(j)-1,\qquad t_p(j),
\]
after discarding indices outside \({1,\ldots,j}\).

Furthermore, for fixed \(p\),
\[
t_p(1)\le t_p(2)\le\cdots\le t_p(n).
\tag{4}
\]

\subparagraph{Proof}

The first assertion follows from the decreasing-then-increasing behavior just described.

For the second, if \(i<t_p(j)\), then
\[
A_i<c(i,j)\le c(i,j+1),
\]
so \(i<t_p(j+1)\). Hence the first crossing cannot move to the left. \(\square\)

Thus, while \(j\) runs from \(1\) to \(n\), a single pointer tracking \(t_p(j)\) moves right at most \(n\) times.

\medskip\hrule\medskip

\subsubsection{7. Precomputing all block costs in \(O(n^2)\)}

By \(1\), the optimal center for ([i,j]) is a point nearest to
\[
m_{ij}=\frac{x_i+x_j}{2}.
\]

For a fixed \(i\), these midpoints are nondecreasing as \(j\) increases. If ties are resolved in favor of the rightmost nearest point, the corresponding optimal-center index is also nondecreasing.

Indeed, between consecutive points \(x_\ell,x_{\ell+1}\), the point \(x_{\ell+1}\) is at least as close as \(x_\ell\) to a number \(m\) exactly when
\[
m\ge \frac{x_\ell+x_{\ell+1}}2.
\]
Once this becomes true, it remains true as \(m\) moves to the right.

Therefore a moving pointer computes all \(c(i,j)\), for fixed \(i\), in \(O(n)\) time. Repeating for all \(i\) gives \(O(n^2)\).

\medskip\hrule\medskip

\subsubsection{8. Exact quadratic-time algorithm}

The following pseudocode uses one-based indices. \texttt{C[i,j]} stores \(c(i,j)\), and \texttt{M[i,j]} stores a corresponding optimal center index.

\begin{verbatim}
QUADRATIC-LINE-K-CENTER(x[1..n], k)

# Phase 1: all one-block costs and centers
for i = 1,...,n:
    ell = i
    for j = i,...,n:
        # Move ell to the rightmost point nearest to (x[i]+x[j])/2.
        # The doubled comparison avoids division by 2.
        while ell < j and
              abs(2*x[ell+1] - x[i] - x[j])
                <= abs(2*x[ell] - x[i] - x[j]):
            ell = ell + 1

        M[i,j] = ell
        C[i,j] = max(x[ell] - x[i], x[j] - x[ell])


# Phase 2: bottleneck dynamic programming
prev[0] = 0
for j = 1,...,n:
    prev[j] = +infinity

for p = 1,...,k:
    cur[0] = 0
    t = 1

    for j = 1,...,n:
        # After this loop, t is the first crossing index;
        # if no crossing exists, t = j+1.
        while t <= j and prev[t-1] < C[t,j]:
            t = t + 1

        candidates = empty set
        if t <= j:
            insert t into candidates
        if t >= 2:
            insert t-1 into candidates

        choose i in candidates minimizing
            max(prev[i-1], C[i,j])

        cur[j] = max(prev[i-1], C[i,j])
        parent[p,j] = i

    prev = cur


# Phase 3: reconstruct the consecutive blocks
F = empty set
p = k
j = n

while j > 0:
    i = parent[p,j]
    insert x[M[i,j]] into F
    j = i - 1
    p = p - 1

# The dynamic program may have used fewer than k nonempty blocks.
while |F| < k:
    insert any point of X \ F into F

return F and prev[n]
\end{verbatim}

\paragraph{Theorem 7: correctness and complexity}

The algorithm returns a \(k\)-element set \(F\subseteq X\) satisfying
\[
\delta(X,F)=R_k.
\]
Its running time is
\[
O(n^2+kn),
\]
and its memory requirement is
\[
O(n^2+kn).
\]

\subparagraph{Proof}

The first phase correctly computes every \(c(i,j)\) and a realizing center by Lemma 1 and the monotone-midpoint argument.

For each fixed \(p\), the pointer \texttt{t} is exactly the crossing index from Lemma 6. Since it is nondecreasing in \(j\), it makes at most \(n\) increments. Lemma 6 shows that checking \texttt{t-1} and \texttt{t} evaluates recurrence \(3\) exactly. Thus the dynamic program computes every \(\rho_p(j)\), in particular
\[
\texttt{prev[n]}=\rho_k(n)=R_k.
\]

Following the parent pointers reconstructs a partition realizing this value. The center \texttt{M[i,j]} realizes the one-block cost for each reconstructed block. Hence the resulting set of centers covers \(X\) with radius \(R_k\).

If fewer than \(k\) centers occur, adding arbitrary unused points cannot increase the radius. It cannot decrease it below \(R_k\), since \(R_k\) was already the minimum over all sets of at most \(k\) centers. Thus the final \(k\)-element set remains optimal.

The cost table takes \(O(n^2)\) time and space. For every \(p\), the dynamic-programming scan takes \(O(n)\), giving \(O(kn)\). The parent table uses \(O(kn)\) space. \(\square\)

\medskip\hrule\medskip

\subsubsection{9. A simpler independent decision algorithm}

For a proposed radius \(r\), one can decide in linear time whether \(k\) centers suffice.

Take the leftmost uncovered point \(x_i\). Choose as center the rightmost \(x_j\) satisfying
\[
x_j\le x_i+r.
\]
This center covers all points through \(x_j+r\). Repeat with the next uncovered point.

This greedy rule uses the minimum possible number of radius-\(r\) centers. Indeed, every feasible cover must have some center \(x_\ell\le x_i+r\) covering \(x_i\). The greedy center satisfies \(x_j\ge x_\ell\), so
\[
x_\ell+r\le x_j+r;
\]
the competing center cannot cover any point farther to the right than the greedy center. Induction on the uncovered suffix proves optimality.

Since an optimal radius is necessarily one of the pairwise distances
\[
x_j-x_i,
\]
sorting all \(O(n^2)\) candidate distances and binary-searching with this decision procedure gives a simpler \(O(n^2\log n)\)-time exact algorithm. The preceding dynamic program removes the logarithmic factor.

\medskip\hrule\medskip

\subsection{Part \(b\): an arbitrary finite metric space}

Let ((X,d)) be a finite metric space with \(|X|=n\), and write
\[
d(x,F)=\min_{f\in F}d(x,f).
\]

\subsubsection{10. Threshold-graph characterization}

For \(r\ge0\), define a graph \(G_r\) on vertex set \(X\) by declaring distinct (u,v) adjacent when
\[
d(u,v)\le r.
\]
A subset \(F\subseteq X\) satisfies
\[
\delta(X,F)\le r
\]
if and only if every vertex is either in \(F\) or adjacent in \(G_r\) to a vertex of \(F\). In other words, \(F\) is a dominating set of \(G_r\).

Let \(\gamma(G)\) denote the minimum size of a dominating set of \(G\), and let
\[
D={d(x,y):x,y\in X}.
\]
Because \(X\) is finite, \(\delta(X,F)\) is always an element of \(D\). Therefore:

\paragraph{Proposition 8}

The exact optimal radius is
\[
\boxed{
R_k^{\mathrm{met}}
=
%
\min{r\in D:\gamma(G_r)\le k}.
}
\tag{5}
\]

Moreover, the optimal \(k\)-element sets are precisely the \(k\)-element dominating sets of \(G_{R_k^{\mathrm{met}}}\).

This is a complete mathematical characterization, but finding a minimum dominating set is computationally difficult.

\medskip\hrule\medskip

\subsubsection{11. An exact finite algorithm}

Assuming the distance matrix is available, one may simply enumerate all \(k\)-element subsets:

\begin{verbatim}
EXACT-METRIC-K-CENTER(X, d, k)

best_radius = +infinity
best_set = none

for every k-element subset F of X:
    r = 0
    for every x in X:
        r = max(r, min_{f in F} d(x,f))

    if r < best_radius:
        best_radius = r
        best_set = F

return best_set and best_radius
\end{verbatim}

There are \(\binom nk\) subsets, and evaluating one requires \(O(nk)\) distance comparisons. Thus the running time is
\[
O\!\left(\binom nk nk\right).
\tag{6}
\]
For every fixed \(k\), this is polynomial in \(n\), but it is exponential when \(k\) is part of the input.

\medskip\hrule\medskip

\subsubsection{12. NP-completeness of the general metric problem}

Consider the decision problem:

Given a finite metric ((X,d)), an integer \(k\), and a radius \(r\), does there exist \(F\subseteq X\), \(|F|\le k\), such that \(\delta(X,F)\le r\)?

It belongs to NP: a proposed \(F\) can be checked in \(O(nk)\) time.

We prove NP-hardness using the standard NP-completeness of \textbf{VERTEX COVER}, one of Karp's classical NP-complete problems. ([Informatique Cornell Bowers][1])

\paragraph{Reduction}

Let \(G=(V,E)\) be a nontrivial instance of VERTEX COVER, with integer \(k\). Construct a graph \(H\) whose vertex set is the disjoint union
\[
W=V\sqcup E,
\]
where each edge \(e\in E\) is regarded as a new ``edge-vertex.''

The adjacencies of \(H\) are:

\begin{enumerate}
\item all vertices of \(V\) are pairwise adjacent, so \(V\) induces a clique;
\item an edge-vertex \(e={u,v}\) is adjacent exactly to \(u\) and \(v\);
\item distinct edge-vertices are not adjacent.
\end{enumerate}

Now define a distance on \(W\) by
\[
d_H(a,b)=
\begin{cases}
0,&a=b,\\
1,&a\ne b\text{ and }a,b\text{ are adjacent in }H,\\
2,&\text{otherwise}.
\end{cases}
\tag{7}
\]

\paragraph{Lemma 9}

The function \(d_H\) is a metric.

\subparagraph{Proof}

Symmetry and positive definiteness are immediate. For the triangle inequality, all nonzero distances are at least \(1\), while every distance is at most \(2\). Thus, for three pairwise distinct points,
\[
d_H(a,c)\le2\le d_H(a,b)+d_H(b,c).
\]
Cases with repeated points are immediate. \(\square\)

A set has covering radius at most \(1\) in this metric exactly when it is a dominating set of \(H\).

\paragraph{Lemma 10}

The original graph \(G\) has a vertex cover of size at most \(k\) if and only if \(H\) has a dominating set of size at most \(k\).

\subparagraph{Proof}

Suppose \(C\subseteq V\) is a vertex cover. Every edge-vertex \(e={u,v}\) is adjacent to at least one point of \(C\). Since \(V\) is a clique, every vertex of \(V\) is in \(C\) or adjacent to every member of \(C\). Thus \(C\) dominates \(H\).

Conversely, let \(D\) dominate \(H\). Whenever an edge-vertex
\[
e={u,v}
\]
belongs to \(D\), replace it by either endpoint, say \(u\). Repeating this operation produces a set
\[
C\subseteq V,\qquad |C|\le |D|.
\]

Every edge-vertex \(e\) is incident to a vertex of \(C\):

\begin{itemize}
\item if \(e\in D\), an endpoint was inserted when \(e\) was replaced;
\item if \(e\notin D\), then \(e\) had to be dominated by one of its endpoints in \(D\), since edge-vertices are not adjacent to one another.
\end{itemize}

Therefore \(C\) contains an endpoint of every edge of \(G\), so \(C\) is a vertex cover. \(\square\)

The construction is polynomial. Thus the metric \(k\)-center decision problem is NP-hard even for the fixed threshold \(r=1\) and metrics whose nonzero distances belong to (\{1,2\}). Together with membership in NP:

\paragraph{Theorem 11}

The decision version of the arbitrary-metric problem is NP-complete, even for \(1\)-\(2\) metrics.

Consequently, the exact optimization problem in Q764\(b\) is NP-hard. Unless \(P=NP\), there is no polynomial-time exact algorithm for arbitrary finite metrics when \(k\) is part of the input.

The distinction between ``at most \(k\)'' and ``exactly \(k\)'' is immaterial: a dominating set of size smaller than \(k\) can be augmented by arbitrary vertices.

\medskip\hrule\medskip

\subsubsection{13. A polynomial-time factor-\(2\) algorithm}

Although exact optimization is NP-hard, there is a very simple optimal-threshold approximation.

\paragraph{Farthest-first algorithm}

Choose an arbitrary first center \(c_1\). Having selected
\[
C_t={c_1,\ldots,c_t},
\]
choose
\[
c_{t+1}\in\operatorname*{arg\,max}_{x\in X} d(x,C_t).
\]

\begin{verbatim}
FARTHEST-FIRST(X, d, k)

choose any c[1] in X
F = {c[1]}

for every x in X:
    nearest[x] = d(x, c[1])

for t = 2,...,k:
    choose c[t] maximizing nearest[x]
    insert c[t] into F

    for every x in X:
        nearest[x] = min(nearest[x], d(x,c[t]))

return F
\end{verbatim}

It performs \(O(n)\) work for each new center, hence \(O(nk)\) distance operations.

\paragraph{Theorem 12}

If \(F\) is returned by farthest-first and \(R^*\) is the exact optimal radius, then
\[
\delta(X,F)\le 2R^*.
\]

\subparagraph{Proof}

Let
\[
r=\delta(X,F),
\]
and choose \(c_{k+1}\in X\) satisfying
\[
d(c_{k+1},F)=r.
\]

We claim that the \(k+1\) points
\[
c_1,c_2,\ldots,c_k,c_{k+1}
\]
are pairwise at distance at least \(r\).

Indeed, when \(c_t\) was selected, its distance from the previously selected centers was the then-current maximum covering distance. These maximum distances form a nonincreasing sequence, whose final value is \(r\). Hence for \(s<t\le k\),
\[
d(c_s,c_t)\ge r.
\]
The same assertion holds for \(c_{k+1}\) by its definition.

Now let \(F^*\) be an optimal set of \(k\) centers with radius \(R^*\). Assign each of the \(k+1\) displayed points to an optimal center at distance at most \(R^*\). By the pigeonhole principle, two of them, say \(c_s,c_t\), are assigned to the same optimal center \(f\). The triangle inequality gives
\[
d(c_s,c_t)
\le d(c_s,f)+d(f,c_t)
\le 2R^*.
\]
Since \(d(c_s,c_t)\ge r\), it follows that
\[
r\le2R^*.
\]
\(\square\)

\medskip\hrule\medskip

\subsubsection{14. The factor \(2\) cannot be improved in polynomial time}

The same reduction used for NP-hardness gives a gap instance.

For the \(1\)-\(2\) metric constructed from a VERTEX COVER instance:

\begin{itemize}
\item if a vertex cover of size at most \(k\) exists, then the optimal metric \(k\)-center radius is \(1\);
\item otherwise, no radius-\(1\) solution exists, while every nonempty set of centers covers all points with radius at most \(2\), so the optimum is \(2\).
\end{itemize}

Suppose a polynomial algorithm had approximation factor
\[
2-\varepsilon
\]
for some fixed \(\varepsilon>0\). On an instance of optimum \(1\), it would return a solution of radius strictly smaller than \(2\). Since all possible radii are (0,1,2), its radius would be at most \(1\). On an instance of optimum \(2\), every feasible solution has radius \(2\). The algorithm would therefore decide VERTEX COVER in polynomial time.

Thus:

\paragraph{Theorem 13}

Unless \(P=NP\), no polynomial-time algorithm for arbitrary finite metrics can guarantee an approximation factor strictly smaller than \(2\).

Hence the farthest-first approximation guarantee is best possible in the standard worst-case complexity sense.

\medskip\hrule\medskip

\subsection{15. Boundary and sanity checks}

\paragraph{\(k=n\)}

Taking \(F=X\) gives
\[
\delta(X,X)=0.
\]
The dynamic program returns the partition into singleton blocks.

\paragraph{\(k=1\) on the line}

Only the two extreme points matter. An optimal center is a point of \(X\) nearest to
\[
\frac{x_1+x_n}{2},
\]
and the optimum is
\[
\min_{1\le \ell\le n}
\max{x_\ell-x_1,x_n-x_\ell}.
\]

\paragraph{Nonuniqueness}

For
\[
X={0,2},\qquad k=1,
\]
both (\{0\}) and (\{2\}) are optimal, with radius \(2\).

\paragraph{A two-cluster example}

For
\[
X={0,1,10,11},\qquad k=2,
\]
the optimal consecutive partition is
\[
{0,1}\mid{10,11}.
\]
Each block has discrete one-center radius \(1\), so an optimal answer is
\[
F={1,11},
\qquad
\delta(X,F)=1.
\]

\paragraph{Why the order is decisive}

On the line, nearest-center cells are consecutive intervals, producing the dynamic program. In an arbitrary metric, the radius-\(r\) neighborhoods can encode essentially arbitrary graph domination instances, which is exactly what causes NP-completeness.

\medskip\hrule\medskip

\subsection{Final answer}

For \textbf{Q764\(a\)}, the problem admits an exact quadratic-time solution. Compute the optimal one-center radius \(c(i,j)\) for every consecutive block, then use the bottleneck recurrence
\[
\rho_p(j)
=
%
\min_{1\le i\le j}
\max{\rho_{p-1}(i-1),c(i,j)}.
\]
Monotonicity reduces each dynamic-programming layer to linear time, giving total complexity
\[
O(n^2+kn)=O(n^2).
\]

For \textbf{Q764\(b\)}, the exact problem is NP-complete even for metrics with distances only \(1\) and \(2\). An exact solution can be obtained by enumerating all \(k\)-subsets in
\[
O\!\left(\binom nk nk\right)
\]
time. In polynomial time, farthest-first selection gives the sharp guarantee
\[
\delta(X,F)\le2,\delta_{\mathrm{opt}},
\]
and no factor below \(2\) is possible unless \(P=NP\).

[1]: https://www.cs.cornell.edu/courses/cs722/2000sp/karp.pdf "cs722"
